\clearpage\SourceLeaf{186}%
\pdfbookmark[1]{A New Prospect for Extending the Scope of Astronomy}{aussicht1813}
\fancyhead[L]{\small A NEW PROSPECT FOR ASTRONOMY}
\begin{center}
{\small A NEW PROSPECT\par
FOR EXTENDING THE SCOPE OF ASTRONOMY.\par}
\medskip\hrule\medskip
{\small Neues Göttingisches Taschenbuch zum Nutzen und Vergnügen für das Jahr 1813.\par
Göttingen, published by Heinrich Dieterich.\par}
\medskip\hrule
\end{center}\medskip
As far back as history reaches, the Earth has always been a restless arena of hostile passions. One people displaces another, only to give way to another in turn, sooner or later. One empire rises on the ruins of another. All human creations are destined to pass away, and the generations that the current of time happens to carry past them share in their blessing or curse. This perpetual change in all things on Earth is humbling and disheartening. Yet it lifts our spirits and cheers us to see the human mind steadily advancing in the realm of truth, nowhere more so than in astronomy, the most sublime of the natural sciences.

For two millennia we have watched the beautiful fruit of observation and deep reflection develop and grow amid the storms of the ages, tended by the loving hands of those devoted to astronomy. Each century passes on this precious inheritance for the next to add to. The heavenly one unveils more and more of her marvelous secrets. The stars have been counted. The seemingly immeasurable distance of the Sun and the bodies that circle it has been measured. Their orbits have been determined precisely, and a single, sublimely simple fundamental force fully explains all their intricate motions and all the varied phenomena that arise from them. Measured,
\expandafter\def\csname HistoricalNoteText13\endcsname{The observations of \textsc{Piazzi}, \textsc{Calandrelli} and \textsc{Brinkley} do give a small annual parallax for the bright star in Lyra. But the observations are still open to doubt, and the results do not agree. Moreover, \textsc{Bradley}’s observations yield no appreciable parallax at all for the same star.}%
\expandafter\def\csname HistoricalNoteText14\endcsname{By means of the transits of 1761 and 1769.}%
\expandafter\def\csname HistoricalNoteText15\endcsname{With the aid of the one small inequality in the Moon’s motion that depends on the solar parallax.}%
\SourceLeaf{187}%
indeed weighed, are the Earth, the Sun, the planets and the moons. The speed of a ray of light has been measured. The delusions that once frightened whole peoples have vanished. The heavenly one does not disdain to serve earthly affairs as well. She divides and counts out periods of time for us, surveys and maps countries and continents, and guides our ships safely across the ocean.

Our solar system has been the main stage of astronomy so far. For many centuries to come it will provide inexhaustible material for new investigations. Yet our solar system is far less in the universe than the Earth is in our solar system. It is a drop in the ocean. The host of fixed stars has so far been little more to us than a set of fixed points to which we refer our observations. We have only an inkling of their immense distances. The longest baseline with which we attempt to measure them, the diameter of the Earth’s annual orbit\HistoricalNote{13}, almost disappears by comparison. Yet let us keep hoping to gain some knowledge of these distances. In measuring the distance of the Sun, the small diameter of the Earth gave no direct measurement. Under favorable circumstances observers chose the nearer Venus instead{\renewcommand{\thefootnote}{**)}\HistoricalNote{14}}, and ingenuity later found a way to use the distance of the Moon, which had long been known, as a yardstick for the distance of the Sun{\renewcommand{\thefootnote}{***)}\HistoricalNote{15}}. Perhaps the very brightest of the fixed stars are not the ones nearest to us. What has so far never quite worked with those stars may succeed more readily with an inconspicuous star of much lower magnitude.

There are indeed good reasons for this conjecture. The so-called fixed stars, which include our Sun, are not fixed in the strict sense of the word. Even if they had all once been at rest at the same time, their effects on one another would have made it impossible for them to remain so. There is motion everywhere even in that countless host of fixed stars. Here time as well as space must be measured on a large scale.
\SourceLeaf{188}%
A century brings less change to their positions relative to one another than a day brings to the relative positions of the planets. Yet with the precision of our observations today we can already discern these changes over fifty or a hundred years in a considerable number of stars. Other conditions being equal, the changes will be most noticeable in the \emph{nearest} stars. Curiously enough, they are greater in some fainter stars than in stars of the first magnitude. \emph{Such} stars, more of which may yet be discovered, deserve especially careful observation for their annual parallax, that is, for the way their position changes with the Earth’s position in its orbit.

One star in particular deserves the closest attention in this respect. It is the star of the sixth magnitude in Cygnus, the Swan, which \textsc{Flamstead} designated by the number 61. Of all the stars in which a shift has so far been observed, this one shows the largest, 8 arcminutes and 52 arcseconds in a hundred years. So large a change in its position relative to us implies either very rapid actual motion through space or a comparatively smaller distance. We owe this important observation to Mr.\ \textsc{Bessel} in Königsberg, who was led to it by comparing \textsc{Bradley}’s observations of this star with observations made in the last years of the preceding century.

Another circumstance makes this star’s observed rapid motion much more interesting still. It is actually what is called a \emph{double star}. A good telescope shows that it consists of two stars that appear very close together (they are 20 arcseconds apart and differ little in brightness). The sky contains many such double stars. Until recently astronomers had always held that two stars appear double to us by chance alone, because they lie almost on a single straight line with us, so their distances from us may differ greatly. A few years ago Mr.\ \textsc{Herschel} was the first to challenge this view and argued that the stars in such pairs are actually close together in space and form a system of their own. \textsc{Herschel}’s reasons for this view are admittedly inconclusive. But even on the face of it, it is extremely unlikely
\SourceLeaf{189}%
that chance should have formed so many double stars with so small a separation. It is a much better bet that there is a real connection between them. Our star in the Swan now weighs heavily in favor of \textsc{Herschel}’s view. If the two stars have no further connection with each other, and they change their position in space or we change ours, we can expect with the highest probability that they will soon cease to appear to us as a double star. Small stars often appear in the sky near Jupiter. As the planet moves on, it seems to leave the stars behind. Only the satellites that are actually close to it keep it faithful company. Yet in 50 years our double star has moved across the firmament thirteen times as far as the distance between its two stars, and that distance itself has hardly changed. Only their alignment is less oblique now (in 1755 one was northeast of the other, and now it is almost due east). There can therefore hardly be any doubt that in them we have come to know a new wonder of the structure of the universe, a solar system in which two suns reign in place of one.

In the physical world nothing rules outright. Whatever rules always feels to some degree the effect of what it rules. Those two suns will rule over each other. The second will not orbit the first, and the first will not orbit the second. Both will orbit their common center of mass. The closer they are to each other and the more strongly they act on each other, the sooner they will complete such a mutual revolution. We will be able to judge how far apart they are if our hope that their distance from us can yet be measured proves well founded. Time will tell how long that revolution takes, perhaps only a few centuries. Then we will have taken a great step forward in our knowledge of the starry sky. We will be able to \emph{weigh} those stars, at least to determine their total mass if we cannot weigh them individually. The first step will be followed by others. But who can count the millennia that must first pass before our knowledge of the starry sky can amount to even a shadow of what we have long known about \emph{our} solar system, so vast and so small? And how much there still is to learn even there, always something new.
\begin{center}\rule{0.24\linewidth}{0.4pt}\end{center}
